The nested sampling scheme assumes there exists some prior knowledge on the underlying (free) energy surface to be sampled. The adaptive biasing force method [\cref{fig:sampling_methods}(m)]\cite{darve2008adaptive,rodriguez2004assessing,darve2008adaptive} aims to overcome this problem by utilising the information gathered on-the-fly in an MD run to determine a bias force $F_{bias}$ that should be added to the system, facilitating it to overcome the free energy barriers. Similar to the nested sampling, in ABF, the CVs are first discretised into multiple points $\bm{\zeta}=\{\zeta_1,\zeta_2,\ldots,\zeta_n \}$ along the reaction pathway, where constrained MDs are performed for some time $t$ at each $\zeta_i\in\bm{\zeta}$. For each MD run, the instantaneous force $\bm{F}(\zeta,t)=\langle \partial_{\zeta}\mathcal{H}\rangle_{\zeta^\prime}^{\mbox{\scriptsize{cond.}}}$ is computed, whereby taking the time average over the MD trajectory gives the mean force $\langle \bm{F}(\zeta)\rangle$ that drives the system dynamics. If now one applies a biasing force that is opposite to the mean force, $i.e.$ $\bm{F}_{\mbox{\scriptsize{bias}}}(\zeta)=-\langle \bm{F}(\zeta)\rangle$ for the subsequent MD simulation, it will help to drive the system to cross the energy barriers by overcoming the mean forces. In this case, \cref{eq:thermodynamic_integration} for free energy computation is modified to 
\begin{equation*}
    W=-\int \bm{F}_{\mbox{\scriptsize{bias}}}(\zeta)\mbox{d}\zeta.
\end{equation*}